% formal/combinatorial-boundary-regularity.tex — boundary-characterization experiment
%
% Defines combinatorial type, boundary events, step-bound formulas,
% and proves continuity of sys across combinatorial boundaries.
% Also contains the gradient formula used by boundary-characterization and multiple-crossings experiments.
%
% Sources:
%   - main.rs: compute_step_bound_detailed(), compute_sys_gradient_a()
%   - experiments/combinatorial-cells/boundary-characterization/: empirical continuity observation
%   - formal/ehz-kkt-system.tex: transition feasibility [lem:numerical-transition-feasibility]
%   - formal/symplectic-polytope-geometry.tex: EHZ capacity [def:ehz-capacity], systolic ratio [def:systolic-ratio]

\subsection{Combinatorial type and boundaries}
\label{sec:combinatorial-boundaries}

\begin{definition}[Combinatorial type]\label{def:combinatorial-type}
Let $K = \{x \in \mathbb{R}^4 : a_k^T x \le 1,\; k = 1, \ldots, F\}$ be a
convex polytope with dual vertices $a_k$.
The \emph{combinatorial type} of~$K$ is the pair $(\mathcal{I}, \Omega)$ where:
\begin{itemize}
  \item $\mathcal{I} \subseteq \{1, \ldots, V\} \times \{1, \ldots, F\}$ is the
    vertex-facet incidence (vertex~$v$ is incident to facet~$k$ iff $a_k^T v = 1$);
  \item $\Omega \in \{-1, 0, +1\}^{F \times F}$ is the $\omega_0$ sign matrix,
    $\Omega_{ij} = \operatorname{sign}(\omega_0(a_i, a_j))$.
\end{itemize}
\end{definition}

\begin{definition}[Transition matrix]\label{def:transition-matrix}
The \emph{transition matrix} $T \in \{0,1\}^{F \times F}$ has $T_{ij} = 1$ iff
facets $i$ and $j$ share at least one vertex \emph{and} $\omega_0(a_i, a_j) \ge 0$.
A cyclic permutation $\sigma = (k_1, \ldots, k_m)$ is \emph{feasible} iff
$T_{k_\ell, k_{\ell+1}} = 1$ for all $\ell$ (indices mod~$m$);
see Lemma~\ref{lem:numerical-transition-feasibility}.
\end{definition}

\begin{definition}[Boundary events]\label{def:boundary-events}
Along a path $a_k(t) = a_k + t \, d_k$ in dual-vertex space, the combinatorial
type changes at the smallest positive~$t$ where one of the following occurs:
\begin{enumerate}
  \item \textbf{Incidence flip:} a vertex's slack w.r.t.\ a non-incident facet
    reaches zero.
  \item \textbf{$\omega_0$ flip:} $\operatorname{sign}(\omega_0(a_i(t), a_j(t)))$
    changes for a pair $(i,j)$ with $\omega_0(a_i, a_j) \ne 0$.
  \item \textbf{Dual vertex degeneration:} $|a_k(t)| \to 0$.
\end{enumerate}
\end{definition}

\begin{lemma}[Incidence flip detection]\label{lem:step-bound-incidence}
Let $v$ be a simple vertex determined by facets $\{j_1, \ldots, j_4\}$, so
$A_v \, v = \mathbf{1}$ where $(A_v)_i = a_{j_i}^T$.
Along $a_k(t) = a_k + t \, d_k$, the vertex moves as
$\dot{v} = -A_v^{-1} \, r$ where $r_i = d_{j_i}^T v$.
For non-incident facet~$j$, the slack $s_j(t) = 1 - a_j(t)^T v(t)$ satisfies
$\dot{s}_j = -d_j^T v - a_j^T \dot{v}$ (linearized).
The incidence flip time is $t_{\mathrm{flip}} = s_j(0) / (-\dot{s}_j)$
when $\dot{s}_j < 0$.
\end{lemma}
\begin{proof}
Differentiate $A_v(t) \, v(t) = \mathbf{1}$ at $t = 0$.
The $i$-th row gives $d_{j_i}^T v + a_{j_i}^T \dot{v} = 0$,
so $A_v \, \dot{v} = -r$ and $\dot{v} = -A_v^{-1} r$.
Then $\dot{s}_j = -(d_j^T v + a_j^T \dot{v})$.
The linearization is exact at first order; the critical time is
$t_{\mathrm{flip}} = s_j(0) / |\dot{s}_j|$ for the closest flip.
% [GAP - This is a first-order approximation; the actual flip time may differ
% slightly due to higher-order terms. The code uses this linear estimate.]
\end{proof}

\begin{lemma}[$\omega_0$ flip detection]\label{lem:step-bound-omega}
Along $a_k(t) = a_k + t \, d_k$, the quantity
$\omega_0(a_i(t), a_j(t))$ is a quadratic polynomial in~$t$:
\[
  \omega_0(a_i(t), a_j(t))
  = \underbrace{\omega_0(a_i, a_j)}_{c}
  + \underbrace{\bigl(\omega_0(d_i, a_j) + \omega_0(a_i, d_j)\bigr)}_{b} \, t
  + \underbrace{\omega_0(d_i, d_j)}_{a} \, t^2.
\]
The sign flips at the smallest positive root of $a t^2 + b t + c = 0$.
\end{lemma}
\begin{proof}
By bilinearity of $\omega_0$:
$\omega_0(a_i + t d_i, a_j + t d_j)
 = \omega_0(a_i, a_j) + t\,[\omega_0(d_i, a_j) + \omega_0(a_i, d_j)]
   + t^2\,\omega_0(d_i, d_j)$.
This is exact (no approximation).
\end{proof}

\subsection{Continuity of the systolic ratio}
\label{sec:sys-continuity}

% For thesis use, prefer Approach 1 below: Hausdorff continuity of the EHZ
% capacity. The direct polytope-specific mechanism is retained only to explain
% why the tempting finite-combinatorics argument does not by itself prove full
% continuity. Local source pointer for wording: papers/ch2021/s1_introduction_and_main_results.tex.

\begin{proposition}[Continuity of $\operatorname{sys}$ at combinatorial boundaries]%
\label{prop:sys-continuous}
Let $a(t) = a + t \, d$ be a path in dual-vertex space crossing a combinatorial
boundary at $t = t^*$. Assume the polytope $K(t)$ is well-defined (bounded, full-dimensional)
for $t$ in a neighborhood of $t^*$. Then
$\operatorname{sys}(K(t))$ is continuous at $t = t^*$.
\end{proposition}
\begin{proof}[Proof sketch]
Write $\operatorname{sys} = c_{\mathrm{EHZ}}^2 / (2V)$.
The volume $V(K(t))$ depends continuously on the dual vertices (it is continuous
in the Hausdorff topology on convex bodies, and $t \mapsto K(t)$ is Hausdorff-continuous
since the half-space description varies linearly).
It suffices to show $c_{\mathrm{EHZ}}(K(t))$ is continuous.

\textbf{Approach 1 (general theory).}
The EHZ capacity is continuous on the space of convex bodies with the Hausdorff
metric~\cite{HoferZehnder1994}. Since $t \mapsto K(t)$ is Hausdorff-continuous,
$c_{\mathrm{EHZ}}(K(t))$ is continuous.
% [Citation note: \cite{HoferZehnder1994}, Ch. 2 (monotonicity axiom A1).
%   No numbered theorem — continuity follows from: K_n → K implies c(K_n) → c(K) via squeezing.
%   Polytopes covered directly by the axiom (no approximation needed).
%   CH2021 states this extension explicitly in the introduction; see
%   papers/ch2021/s1_introduction_and_main_results.tex.
%   For explicit discussion: Cieliebak-Hofer-Latschev-Schlenk (2007), MSRI Publ. vol. 54.
%   Additional refs: McDuff--Salamon (1998), Exercise~12.7 (cited in HK2017 proof of Thm~1.8);
%   Artstein-Avidan \& Ostrover (2014), Prop~2.7 for non-smooth extension to polytopes.]

\textbf{Approach 2 (polytope-specific mechanism).}
For a polytope $K$, write $c_{\mathrm{EHZ}}(K) = \min_{\sigma \in \mathcal{F}(K)} A(\sigma)$,
where $\mathcal{F}(K)$ is the set of feasible cyclic permutations
(those forming feasible cycles in the transition matrix,
Definition~\ref{def:transition-matrix})
and $A(\sigma) = 1/(2\,Q(\sigma))$ is the action.
Here $Q(\sigma) = \beta^\top H \beta$ where $\beta$ solves the KKT system
for permutation~$\sigma$; it depends continuously on the dual vertices
when the KKT system is non-degenerate.

\emph{Key structural fact:} at a generic codimension-1 combinatorial boundary,
$\mathcal{F}(K(t^*)) \supseteq \mathcal{F}(K(t^* - \delta)) \cup \mathcal{F}(K(t^* + \delta))$
for small $\delta > 0$.
This is because:
\begin{itemize}
  \item At an $\omega_0$ boundary where $\omega_0(a_i, a_j)$ passes through zero:
    at $t^*$, both transitions $i \to j$ and $j \to i$ are feasible ($\omega_0 \ge 0$
    is satisfied with equality). A cyclic permutation visits each facet at most once,
    so it uses at most one of $i \to j$ or $j \to i$, hence is feasible on at least
    one side.
  \item At an incidence boundary: vertex adjacency at $t^*$ is a superset of either side's.
\end{itemize}
There are no ``boundary-only'' orbits (feasible at $t^*$ but on neither side)
at generic codimension-1 boundaries.

\emph{Lower semicontinuity} ($c_{\mathrm{EHZ}}(t^*) \le \liminf$):
Since $\mathcal{F}(t^*) \supseteq \mathcal{F}(t^* \pm \delta)$, the minimum
over the larger set is $\le$ the minimum over each side. So
$c_{\mathrm{EHZ}}(t^*) \le \lim_{t \to t^{*\pm}} c_{\mathrm{EHZ}}(t)$.

\emph{Upper semicontinuity} ($c_{\mathrm{EHZ}}(t^*) \ge \limsup$):
% [GAP - The direct argument does NOT give upper semicontinuity.
% A minimizer $\sigma^*$ at $t^*$ is feasible on at least one side (say the left).
% Then $c(t^* - \delta) \le A(\sigma^*, t^* - \delta) \to A(\sigma^*, t^*) = c(t^*)$.
% So $\limsup_{\text{left}} \le c(t^*)$. But from the right, the minimizer at $t^*$
% might not be feasible, and the right-feasible orbits might have strictly higher
% actions at $t^*$: $\lim_{\text{right}} c(t) = \min_{F(\text{right})} A(\sigma, t^*)
% > c(t^*) = \min_{F(t^*)} A(\sigma, t^*)$.
% This would make c LOWER semicontinuous with a possible downward jump at $t^*$.
% Full continuity requires either the general theory (Approach 1) or an additional
% argument that the minimizer at $t^*$ is feasible on BOTH sides.]
This direction requires Approach~1 (general theory) or an argument that
$c_{\mathrm{EHZ}}(t^*)$ cannot be strictly lower than either one-sided limit.
The polytope-specific mechanism alone gives only lower semicontinuity.
\end{proof}

\begin{remark}\label{rem:sys-continuous-empirical}
The experiment confirms continuity empirically: across 873 boundary crossings,
$\max |\Delta\operatorname{sys}| = 2.91 \times 10^{-4}$
(limited by the finite step-over~$\varepsilon$, not by a discontinuity).
\end{remark}

\subsection{Systolic ratio gradient}
\label{sec:sys-gradient}

\begin{lemma}[Gradient of $\operatorname{sys}$ in dual vertices]\label{lem:sys-gradient-a}
% TODO: add cross-references to capacity_derivatives_a and volume_derivatives_a lemmas
For a polytope $K$ with dual vertices $a_k$, volume $V$, capacity $c$, and
$\operatorname{sys} = c^2 / (2V)$, the gradient with respect to dual vertex~$a_k$ is
\[
  \frac{\partial \operatorname{sys}}{\partial a_k}
  = \frac{1}{V}\!\left(
    c \, \frac{\partial c}{\partial a_k}
    - \operatorname{sys} \, \frac{\partial V}{\partial a_k}
  \right).
\]
\end{lemma}
\begin{proof}
By the quotient rule applied to $\operatorname{sys} = c^2 / (2V)$:
\[
  \frac{\partial \operatorname{sys}}{\partial a_k}
  = \frac{2c \, \frac{\partial c}{\partial a_k} \cdot 2V
    - c^2 \cdot 2\,\frac{\partial V}{\partial a_k}}{(2V)^2}
  = \frac{c \, \frac{\partial c}{\partial a_k}}{V}
    - \frac{c^2}{2V} \cdot \frac{\frac{\partial V}{\partial a_k}}{V}
  = \frac{c \, \frac{\partial c}{\partial a_k} - \operatorname{sys} \, \frac{\partial V}{\partial a_k}}{V}.
  \qedhere
\]
\end{proof}
